\documentclass[10pt,a4paper]{amsart}
\usepackage[margin=19mm]{geometry}
\usepackage{amsmath,amssymb,mathtools}
\usepackage[scr=rsfs,cal=euler]{mathalfa}
\usepackage{xfrac}
\usepackage[unicode,hidelinks,bookmarks=false]{hyperref}
\newcommand{\ie}{i.e.\ }
\newcommand{\cf}{cf.\ }
\newcommand{\resp}{respectively\ }
\newcommand{\Subsection}{\subsection}
\providecommand{\nobreakdash}{\nobreak\hbox{-}}
\setlength{\emergencystretch}{2em}
\multlinegap0pt
\usepackage{amssymb}
\newtheorem{theorem}{Theorem}
\newtheorem{lm}{Lemma}
\newtheorem{prop}{Proposition}
\newcommand{\B}{\mathbf{B}}
\newcommand{\C}{\mathbb{F}}
\newcommand{\GL}{{\mathrm{GL}}}
\newcommand{\Mat}{{\mathrm{Mat}}}
\newcommand{\diag}{\operatorname{diag}}
\newcommand{\g}{\mathfrak g}
\newcommand{\gC}{{\mathfrak g}_{\C}}
\newcommand{\gl}{\mathfrak g \mathfrak l}
\renewcommand{\mid}{\,:\,}
\newcommand{\sgn}{\operatorname{sgn}}
\newcommand{\triv}{\mathrm{triv}}
\renewcommand{\u}{\mathfrak u}
\title{Closed Orbits and Descents for Enhanced Standard Representations of Classical Groups}
\author{Chen Liang}
\date{Source: arXiv:2507.21911v3 (CC BY 4.0)}
\begin{document}
\maketitle

\noindent\emph{Source and licence.} Closed Orbits and Descents for Enhanced Standard Representations of Classical Groups. Version arXiv:2507.21911v3 (CC BY 4.0), distributed by arXiv under the Creative Commons Attribution 4.0 International licence, http://creativecommons.org/licenses/by/4.0/, as recorded in the arXiv licence metadata for this exact version and shown on the abstract page https://arxiv.org/abs/2507.21911v3. Full licence notice: \texttt{licenses/arxiv-2507.21911-CC-BY-4.0.txt}.\par\smallskip

\noindent\textit{Editorial scope.} Counted unit: the article's prop prop:descentglnmvw (source0.tex lines 1725--1733). The article's complete original proof of it is delivered with it (source0.tex lines 1734--1786, one \texttt{proof} environment of 53 lines). No mathematics is written, repaired or re-derived: the statement and the proof are the article's own lines, in the article's own notation, and no retained line is substituted or re-flowed.  Retained uncounted as the source-local context the proof needs: lines 608--611; lines 1707--1722.  Nothing was omitted from any retained span, so the package carries no omission record.  No retained line contains a citation, so the wrapper carries no bibliography and no bibliography entry.  No retained line contains a figure environment, an image-inclusion command or drawing code, so no image is delivered and the package declares no asset dependency.  Editorial formatting only, in addition: an AMS \texttt{amsart} preamble that restates the article's own declarations and macros used by the retained lines, namely the environments \texttt{theorem}, \texttt{lm}, \texttt{prop} on separate counters, together with the standard packages those lines need.  The retained environments print in this document as Theorem 1, Lm 1, Proposition 1 in order of appearance, whereas the article prints the same items with the numbering of its own full document.  The complete native source of the article as distributed in the arXiv e-print for this exact version is bundled unchanged as \texttt{sources/182375/source0.tex}.
    \begin{theorem}
		\label{thm:closedorbitsmvw}
        Every $G$-closed orbit in $\g\times E$ is $\breve{G}$-stable.
	\end{theorem}

    \begin{lm}
    	\label{lm:descentgln}
    	Let $h_1,h_2\in\GL_m(\C)$ be symmetric. For $i=1,2$, consider representations $(\rho_i,\C^{m\times 1}\times\C^{1\times m})$ and $(\omega_i,\gl_m(\C))$ of $\{\pm 1\}$ defined by
    	\begin{equation*}
    		\rho_i(-1)(u,v)=(-h_iv^t,-u^th_i^{-1})\quad\text{and}\quad\omega_i(-1)X=h_iX^th_i^{-1}.
    	\end{equation*}
    	Then
    	\begin{equation*}
    		\rho_1\simeq\rho_2,\quad (u,v)\mapsto(h_2h_1^{-1}u,v)
    	\end{equation*}
    	and
    	\begin{equation*}
    		\omega_1\simeq\omega_2,\quad X\mapsto h_2h_1^{-1}X.
    	\end{equation*}
    	In particular, $\rho_i\simeq\triv^m\oplus\sgn^m$ and $\omega_i\simeq\triv^{m(m+1)/2}\oplus\sgn^{(m-1)m/2}$ for $i=1,2$, where $\sgn$ denotes the sign character.
    \end{lm}

    \begin{prop}
    	\label{prop:descentglnmvw}
    	The stabilizer $\breve{\GL}_n(\C)_x\simeq\breve{\GL}_{n-k}(\C)$, and
    	\begin{equation}
    		\label{eq:descentgln}
    		N_{O,x}^{\g\times E}\simeq(\gl_{n-k}(\C)\times\C^{(n-k)\times 1}\times\C^{1\times(n-k)})\oplus\triv^{2k}
    	\end{equation}
    	as representations of $\breve{\GL}_{n-k}(\C)$.
    \end{prop}

    \begin{proof}
    	We have seen that
    	\begin{equation*}
    		\GL_n(\C)_x=\left\{\diag(I_k,g)\mid g\in\GL_{n-k}(\C)\right\}\simeq\GL_{n-k}(\C).
    	\end{equation*}
    	By Theorem \ref{thm:closedorbitsmvw}, we have $\breve{\GL}_n(\C)x=\GL_{n}(\C)x$, and hence $\GL_n(\C)_x$ is a subgroup of $\breve{\GL}_n(\C)_x$ of index two.
    	
    	Let $y=(y_1,\dots,y_k)$ and $g_0=\diag(-\psi(y),I_{n-k})\in\GL_n(\C)$. Then we have $(g_0,-1)\in\breve{\GL}_n(\C)_x$, and the map $\breve{\GL}_n(\C)_x\to\breve{\GL}_{n-k}(\C)$ defined by
    	\begin{equation*}
    		(\diag(I_k,g),1)\mapsto g\quad\text{and}\quad(g_0,-1)\mapsto(I_{n-k},-1).
    	\end{equation*}
    	is an isomorphism.
    	
    	Now we consider the action of $\breve{\GL}_{n-k}(\C)$ on $N_{O,x}^{\g\times E}$. For $z\in\g\times E$, write
    	\begin{equation*}
    		z=\left(\left(\begin{matrix}
    			A&B\\C&D
    		\end{matrix}\right),\left(\begin{matrix}
    		    u\\u'
    		\end{matrix}\right),(v,v')\right),
    	\end{equation*}
    	where $A\in\gl_k(\C),\ B\in\Mat_{k\times(n-k)}(\C),\ C\in\Mat_{(n-k)\times k}(\C),\ D\in\gl_{n-k}(\C),\ u\in\C^{k\times 1},\ u'\in\C^{(n-k)\times 1},\ v\in\C^{1\times k}$ and $v'\in\C^{1\times (n-k)}$. The action of $\breve{\GL}_{n-k}(\C)$ on $\g\times E$ is given by
    	\begin{equation*}
    		(g,1).z=\left(\left(\begin{matrix}
    			A&Bg^{-1}\\gC&gDg^{-1}
    		\end{matrix}\right),\left(\begin{matrix}
    			u\\gu'
    		\end{matrix}\right),(v,v'g^{-1})\right)
    	\end{equation*}
    	and
    	\begin{equation*}
    		(I_{n-k},-1).z=\left(\left(\begin{matrix}
    			\psi(y)A^t\psi(y)^{-1}&\psi(y)C^t\\B^t\psi(y)^{-1}&D
    		\end{matrix}\right),\left(\begin{matrix}
    			-\psi(y)v^t\\-v'^t
    		\end{matrix}\right),(-u^t\psi(y)^{-1},-u'^t)\right).
    	\end{equation*}
    	The tangent space
    	\begin{equation*}
    		T_xO=\left\{w=\left(\begin{matrix}
    			A&B\\C&0_{(n-k)\times(n-k)}
    		\end{matrix}\right)\mid A\in\gl_k(\C),\ B\in\Mat_{k\times(n-k)}(\C),\text{ and }C\in\Mat_{(n-k)\times k}(\C)\right\}
    	\end{equation*}
    	with the action of $\breve{\GL}_{n-k}(\C)$ given by
    	\begin{equation*}
    		(g,1).w=\left(\begin{matrix}
    			A&Bg^{-1}\\gC&0_{(n-k)\times(n-k)}
    		\end{matrix}\right)\quad\text{and}\quad(I_{n-k},-1).w=\left(\begin{matrix}
    		-\psi(y)A^t\psi(y)^{-1}&\psi(y)C^t\\B^t\psi(y)^{-1}&0_{(n-k)\times(n-k)}
    		\end{matrix}\right).
    	\end{equation*}
    	Comparing these two representations, we obtain \eqref{eq:descentgln} by Lemma \ref{lm:descentgln}.
    \end{proof}

\end{document}
